\chapter{Lições Aprendidas}

%\textcolor{red}{Documentar as lições aprendidas de modo a aperfeiçoar os processos e evitar que os erros e problemas encontrados se repitam em futuros projetos.}

\section{Análise SWOT}

A análise SWOT (\textit{Strengths}, \textit{Weaknesses}, \textit{Opportunities}, \textit{Threats}) é usada para identificar os pontos fortes e fracos do produto, e as principais oportunidades e ameaças normalmente associadas aos pontos fortes e fracos identificados. A Tabela \ref{tab:swot} apresenta a análise SWOT do micromouse desenvolvido.

\begin{table}[ht]
% \label{tab:swot}
\centering
\caption{Análise SWOT do produto.}
\begin{tabular}{|c|c|c|} \hline
 & \textbf{Pontos fortes} & \textbf{Pontos fracos}  \\ \hline 
    {\parbox{0.1\textwidth}{ \textbf{Fatores \\ internos}}}  &
    {\parbox{0.4\textwidth}{ Forças:
        \begin{itemize}
            \item Redução de custo com utilização de alguns componentes nacionais.
            \item Telemetria em tempo real com 90,66\% de cobertura de testes.
            \item Chassi utiliza material leve, resistente e sob medida com impressão 3D.
        \end{itemize} }} &
    {\parbox{0.4\textwidth}{ Fraquezas:
        \begin{itemize}
            \item Atrasos na impressão do chassi, por causa de problemas na impressão 3D.
            \item Dificuldade para encontrar peças de reposição localmente.
            \item Retrabalho por troca de componentes já adquiridos.
        \end{itemize} }}\\ \hline
    {\parbox{0.1\textwidth}{ \textbf{Fatores \\ externos}}}  &
    {\parbox{0.4\textwidth}{ Oportunidades:
        \begin{itemize}
            \item Mercado de robótica educacional em expansão, com poucas soluções nacionais de baixo custo.
            \item Possível participação em competições oficiais IEEE \textit{Micromouse}.
            \item Reaproveitamento da plataforma em disciplinas da FCTE.
        \end{itemize} }} &
    {\parbox{0.4\textwidth}{ Ameaças:
        \begin{itemize}
            \item Dependência de componentes importados, sujeitos a atraso e câmbio.
            \item Exigências ANATEL em caso de comercialização.
            \item Concorrência de equipes internacionais mais experientes.
        \end{itemize} }}\\ \hline
\end{tabular}
\label{tab:swot}
\end{table}


\section{Planejado x realizado}
\subsection{Objetivos}
%\textcolor{red}{Os objetivos foram alcançados? Responda as questões e comente os pontos mais relevantes – até 700 caracteres, considerando os espaços.}

O chassi proposto inicialmente foi utilizado, mas não da maneira planejada: o plano original previa uma placa perfurada de fibra de vidro com os componentes eletrônicos soldados, e acabou-se utilizando duas protoboards coladas sobre o chassi. A bateria alimentou o robô tranquilamente durante todo o tempo em que ele ficou ligado, e os componentes eletrônicos propostos foram suficientes para que o micromouse atendesse aos requisitos necessários. A solução de software proposta para a telemetria foi alcançada, com as informações coletadas pelo robô aparecendo corretamente no \textit{frontend}. Entretanto, a integração entre os sensores, os motores e o algoritmo de exploração do labirinto não funcionou, devido ao pouquíssimo tempo de teste disponível com o robô montado, de modo que o robô não conseguiu concluir nenhum dos labirintos.

\subsection{Prazos}
%\textcolor{red}{O projeto foi entregue dentro do prazo? Responda as questões e comente os pontos mais relevantes – até 700 caracteres, considerando os espaços.}

Houve demora na montagem do micromouse devido ao atraso na entrega de componentes eletrônicos e de peças do chassi, a um erro na impressão 3D do chassi que levou dias para ser solucionado e à necessidade de remodelar o chassi para reduzir suas dimensões, já que a versão anterior estava muito grande. Mesmo após o chassi estar pronto, a soldagem dos componentes eletrônicos na placa também demorou, e, com o micromouse totalmente montado, descobriu-se que o Raspberry Pi Pico W estava queimado. Foi então necessário dessoldar todos os componentes, substituí-los pela ESP32 WROOM DevKit V1, remontá-los em uma protoboard e adaptar o firmware, originalmente escrito para a Pico W, à nova placa, o que fez com que o robô não estivesse pronto para completar um labirinto a tempo da apresentação.

\subsection{Orçamento}
%\textcolor{red}{O projeto foi entregue dentro do orçamento? Responda as questões e comente os pontos mais relevantes – até 700 caracteres, considerando os espaços.}

O grupo acabou gastando mais do que o orçamento inicial havia previsto, uma vez que foram necessárias diversas compras adicionais de componentes: uma bateria recebida veio com defeito e um dos motores não funcionava, exigindo a aquisição de novas unidades. Ao final do projeto, o orçamento total foi atualizado para refletir essas compras extras.

\subsection{Escopo}
%\textcolor{red}{O  projeto atendeu ao escopo? Responda as questões e comente os pontos mais relevantes – até 700 caracteres, considerando os espaços.}

O escopo completo do projeto não foi atingido, já que o robô não conseguiu completar o labirinto. Entretanto, o micromouse era capaz de se locomover de forma autônoma e de enviar dados de telemetria para o \textit{frontend}.


\section{Projeto}
%\textcolor{red}{Comente os pontos mais relevantes a serem aperfeiçoados ou adotados em próximos projetos.}
\subsection{Pontos fortes}
\begin{enumerate}
	\item Resiliência da equipe para contornar problemas técnicos que surgiam constantemente ao longo do projeto.
	\item Capacidade de adaptação a necessidades que surgiam no meio do desenvolvimento do projeto.
	\item Aprendizado sobre como realizar uma boa pesquisa de mercado na compra de componentes eletrônicos e de peças do chassi.
	%\item \textcolor{red}{Até 200 caracteres, considerando os espaços.}
\end{enumerate}
\subsection{Pontos fracos}
\begin{enumerate}
	\item Alta dependência dos líderes das áreas de Eletrônica e de Estrutura na elaboração do chassi e na montagem dos componentes eletrônicos.
	\item Demora na montagem final do robô, o que impediu a realização de testes de integração a tempo da entrega final.
	%\item \textcolor{red}{Até 200 caracteres, considerando os espaços.}
	%\item \textcolor{red}{Até 200 caracteres, considerando os espaços.}
\end{enumerate}

\section{Recomendações para projetos futuros}

\begin{enumerate}
	\item Comprar materiais em lojas físicas, o que facilita encontrar um substituto rapidamente caso seja necessário adquirir uma nova unidade.
	\item Comprar o dobro da quantidade necessária de cada componente, por precaução.
	\item Ter um cronograma mais rigoroso.
	%\item \textcolor{red}{Até 200 caracteres, considerando os espaços.}
\end{enumerate}

\section{Questões em aberto}

%\textcolor{red}{Usar caso haja alguma questão pendente em relação às entregas do projeto (Ex.: Requisitos não entregues, melhoria na programação do software, troca de sensores, entre outros.}

\begin{enumerate}
	\item O robô não conseguiu completar o labirinto, de modo que não houve uma integração completa de todos os componentes eletrônicos.
	%\item \textcolor{red}{Até 200 caracteres, considerando os espaços.}
	%\item \textcolor{red}{Até 200 caracteres, considerando os espaços.}
	%\item \textcolor{red}{Até 200 caracteres, considerando os espaços.}
\end{enumerate}

%\section{Desempenho dos fornecedores}
%
%\subsection{Fornecedores com desempenho acima do esperado}
%\textcolor{red}{Fornecedores que você certamente convidaria para um próximo projeto. Citar AAAA – justificativa. Até 200 caracteres, considerando os espaços.}
%
%\subsection{Fornecedores com desempenho conforme esperado}
%\textcolor{red}{Fornecedores que você provavelmente convidaria para um próximo projeto. Citar AAAA – justificativa. Até 200 caracteres, considerando os espaços.}
%
%\subsection{Fornecedores com desempenho abaixo do esperado}
%\textcolor{red}{Fornecedores que você não convidaria para um próximo projeto. Citar AAAA – justificativa. Até 200 caracteres, considerando os espaços.}